\documentclass[aspectratio=169,10pt,spanish]{beamer}

\input{../../../_preambulo_beamer.tex}
\input{../../../_comandos_beamer.tex}

\selectlanguage{spanish}

\title{MA1001B - Modelación Estadística para la Toma de Decisiones}
\subtitle{Lección 16: Distribución hipergeométrica: muestreo sin reemplazo}
\author{}
\institute{Tecnol\'ogico de Monterrey}
\date{}

\begin{document}

\begin{frame}[plain]
  \vspace{0.9cm}
  {\Large\bfseries MA1001B - Modelación Estadística para la Toma de Decisiones\par}
  \vspace{0.5cm}
  {\large Lección 16: Distribución hipergeométrica\par}
  \vspace{0.1cm}
  {\large Muestreo sin reemplazo\par}
  \vspace{0.5cm}
  {\color{TecRojo}\rule{\textwidth}{1.2pt}}
  \vspace{0.6cm}

  Facultad MA1001B\\[0.3cm]
  Periodo\\[0.3cm]
  Tecnol\'ogico de Monterrey
\end{frame}

\begin{frame}{El lote es finito}
  \begin{alertblock}{Pregunta guía}
    Si $6$ de $80$ unidades son defectuosas y se extraen $10$ sin reemplazo,
    ¿cuál es $P(X=0)$ y por qué una aproximación binomial es demasiado
    optimista?
  \end{alertblock}

  \vspace{0.6em}
  Al final de esta lección, usted puede:
  \begin{itemize}
    \item identificar $N$, $K$ y $n$ en una extracción sin reemplazo;
    \item calcular $P(X=0)$ hipergeométrica;
    \item calcular $E[X]=nK/N$;
    \item explicar por qué la binomial sobreestima $P(X=0)$ cuando $n/N$ no es pequeña.
  \end{itemize}

  \vfill
  \scriptsize Todos los lotes usados hoy son sintéticos. No se usan datos reales de empresa.
\end{frame}

\begin{frame}{Ingredientes hipergeométricos}
  \small
  \[
    N=80,\qquad K=6,\qquad n=10,\qquad \frac{n}{N}=0.125.
  \]
  $X$ es el número de defectuosos en la extracción. Como las unidades no se
  reponen, cada extracción cambia el lote restante, de modo que los ensayos
  no son independientes.
\end{frame}

\begin{frame}[fragile]{Paso 1: $P(X=0)$ sin reemplazo}
  \begin{columns}[T]
    \begin{column}{0.42\textwidth}
      \small
      \begin{lstlisting}
p0 = stats.hypergeom.pmf(
    0, M=80, n=6, N=10
)
      \end{lstlisting}

      \begin{block}{Salida verificada}
        $n/N=0.125000$\\
        $P(X=0)=0.436326$
      \end{block}
    \end{column}
    \begin{column}{0.58\textwidth}
      \centering
      \includegraphics[width=\textwidth]{../figuras/hipergeometrica_01_sin_reemplazo.png}
      \vspace{0.2em}
      \scriptsize FMP hipergeométrica del Paso 1
    \end{column}
  \end{columns}
\end{frame}

\begin{frame}{La media no necesita independencia}
  \small
  \[
    E[X]=\frac{nK}{N}=\frac{10\times 6}{80}=0.75.
  \]
  Esa media coincide con una binomial con $p=K/N=0.075$. Coincidir en la
  media no coincide en $P(X=0)$.
\end{frame}

\begin{frame}[fragile]{Paso 2: Esperanza hipergeométrica}
  \begin{columns}[T]
    \begin{column}{0.40\textwidth}
      \small
      \begin{lstlisting}
mean = n * K / N
      \end{lstlisting}

      \begin{block}{Salida verificada}
        $nK/N=0.750000$\\
        media scipy $=0.750000$
      \end{block}
    \end{column}
    \begin{column}{0.60\textwidth}
      \centering
      \includegraphics[width=\textwidth]{../figuras/hipergeometrica_02_esperanza.png}
      \vspace{0.2em}
      \scriptsize Media marcada en la FMP del Paso 2
    \end{column}
  \end{columns}
\end{frame}

\begin{frame}{Paso 3: La binomial sobreestima $P(X=0)$}
  \begin{columns}[T]
    \begin{column}{0.42\textwidth}
      \small
      Aproximación binomial: $X\sim\mathrm{Binomial}(10,0.075)$.
      \[
        P_{\text{binom}}(X=0)=0.458582
      \]
      frente a hipergeométrica $0.436326$. Sobreestimación $0.022257$.

      \vspace{0.4em}
      \textbf{Límite:} la binomial sobreestima $P(X=0)$ cuando $n/N$ no es
      pequeña. La Lección 17 modela eventos raros con una tasa de Poisson.
    \end{column}
    \begin{column}{0.58\textwidth}
      \centering
      \includegraphics[width=\textwidth]{../figuras/hipergeometrica_03_aproximacion_binomial.png}
      \vspace{0.2em}
      \scriptsize Hipergeométrica frente a binomial del Paso 3
    \end{column}
  \end{columns}
\end{frame}

\begin{frame}{Verificación de conceptos}
  \begin{enumerate}
    \item ¿Por qué las diez extracciones no son independientes?
    \item Calcule $E[X]=nK/N$.
    \item ¿Cuál modelo da la $P(X=0)$ más grande?
    \item ¿Por qué $n/N=0.125$ es una advertencia para la aproximación binomial?
  \end{enumerate}

  \vspace{0.6em}
  \begin{block}{Conexión con la Lección 17}
    La sesión puede continuar directamente con el modelo de Poisson para
    eventos raros en un turno.
  \end{block}

  \vfill
  \scriptsize Fuente: OpenStax, \textit{Introductory Business Statistics 2e},
  Sección 4.1. CC BY-NC-SA.
\end{frame}

\appendix



\end{document}
